\section{局限与影响}
\label{sec:limitations}

当前绑定约束是损伤分类器在事件不相交测试上对中间损伤等级（轻微/严重损伤）召回持续偏低、开放词汇指认存在误报，而不是仿真平台或导航几何
（oracle 指认下、指认仍生效时到达 SR $=1.00$，GPS 扰动 $\sigma\le 10\,\mathrm{m}$ 对 SR 无单调影响；该结论以本仿真的运动学假设——无风扰、无避障、无能耗约束——为前提，不外推到真实飞行）。
类别加权训练把训练集 macro-F1 提升到 $0.49$--$0.57$，证明分类器有能力学到损伤线索，但同等条件下验证/测试 macro-F1 仅到 $0.22$--$0.28$，
是真实的事件级泛化缺口，不是训练配方完全未处理不平衡的产物。
与标准（非事件不相交）切分对照显示：换用文献常见的更宽松协议后 macro-F1 仅升至 $0.31$，
离公开基线的 $\sim 0.7$ 仍有明显差距，这部分差距应归因于本文使用的轻量架构（ResNet18 双塔、8 epoch、无集成）与训练规模，
更强的骨干网络、更长训练与专门的类别不平衡技巧（focal loss、两阶段训练、过采样）留作后续工作。
分辨率阶梯会改变高频像素，降质后平均熵方向已与假设一致（随 GSD 变粗上升），但幅度仍小，不足以支撑稳健的跨 GSD 触发决策。
该阶梯是在单一原生分辨率上撤销合成模糊，零结果只说明机制在该构造上不可识别，不能外推为真实低空厘米级观测无信息。
RescueNet 提供真实低空端点，但没有与 xBD 瓦片的逐栋地理配准，且缺乏灾前配对（差通道恒为 0），该结果只能作为域外可用性检查，不能单独支持``绑定约束在损伤头''的因果论证。
语言条件搜索的指令空间仍偏模板化，尽管已加入方位、否定与多地标；在语言复杂度不足时不应自称完整 VLN。
旧 40$\times$3 设计统计功效不足；新基准规模按功效分析扩大，但仍需损伤证据阳性回合才能比较复检策略。
仿真不包含风扰、斜视、运动模糊、通信与电池。
\ProjectName 定位为决策支持，损伤结论应保留人工确认。
地理影像可能暴露敏感位置，轨迹发布需按灾种做隐私审查。
